\subsection{A Motivating Example}

\begin{frame}{\insertsubsection}
	\slideWaterfallModel
\end{frame}

\begin{frame}{\insertsubsection}
	\begin{fancycolumns}[animation=none]
		\begin{tcolorbox}[colback=red!25!white,colframe=red!70!black,title=Bus]
			As a citizen I would like to take the bus to move around the city so that I can travel without having to drive.
		\end{tcolorbox}\pause
		
		\begin{tcolorbox}[colback=red!25!white,colframe=red!70!black,title=Bus Stop]
			As a citizen I want to have a simple bus stop so I can get on the bus close to where I life.
		\end{tcolorbox} \pause
		
		\Huge\dots \pause
		\nextcolumn
		\begin{note}{Targets}
			\begin{itemize}
				\item Design Specification
				\item Source Code
				\item Test Code
				\item Code and Test Documentation
				\item \dots
			\end{itemize}
		\end{note}\pause
		\begin{note}{Desired Properties}
			\begin{itemize}
				\item Human-independent
				\item Correct
				\item Maintainable
				\item Effective
				\item \dots
			\end{itemize}
		\end{note}
	\end{fancycolumns}
\end{frame}

\subsection{Model-driven Software Engineering}
\begin{frame}{\insertsubsection{} \mytitlesource{\href{https://link.springer.com/book/10.1007/978-3-031-02549-5}{MDSE}}}
	\begin{fancycolumns}[animation=none]
		\begin{definition}{Model-driven Software Engineering}
			Software development methodology where domain models are the primary artifact
		\end{definition} \pause
		\begin{note}{Central Process}
			\begin{enumerate}
				\item Domain experts help create domain model
				\item Domain model is evaluated/transformed
				\item Other artifacts are (semi-)automatically created
			\end{enumerate}
			
		\end{note}\pause
		\begin{note}{Advantages}
			\begin{itemize}
				\item Model Checking
				\item Model-to-Model Transformations
				\item Model-to-Code Transformations
				\item \dots
			\end{itemize}
		\end{note}
		\nextcolumn
		\begin{note}{UML Modelling Layers}
			\begin{enumerate}
				\item Originals (Objects in memory)
				\item Models (Object diagram)
				\item Meta-Models (Class diagram)
				\item Meta-Meta-Models (Diagram diagram)
			\end{enumerate}
		\end{note}
		\begin{example}{Exemplary EMF Meta-Model}
			\centering\pic[width=\columnwidth]{ml/emf}
		\end{example}
		
	\end{fancycolumns}
\end{frame}

\subsection{Model-to-Code Transformation}
\begin{frame}{\insertsubsection{}}
	\begin{fancycolumns}[animation=none]
		\begin{definition}{Model-to-Code Transformation}
			\begin{itemize}
				\item Special case of Model-to-Model transformation
				\item Typically at the end of a chain of model transformations
				\item Typically not all model information $\Rightarrow$ information loss
				\item Not necessarily code (e.g., XML files)
			\end{itemize}
		\end{definition}\pause
		\begin{note}{Motivation}
			\begin{itemize}
				\item Manual code generation takes effort
				\item Manual code generation is error-prune
				\item Model-code-divergence possible
			\end{itemize}
		\end{note} \pause
		\nextcolumn
		\begin{note}{Code Generators}
			\begin{itemize}
				\item Template-based code generation
				\item Static parts mixed with dynamic section
				\item Dynamic sections filled using queries (comparable to SQL)
				\item Custom-defined for a specific model
			\end{itemize}
		\end{note} \pause
		\begin{example}{Xtend}
			\centering\pic[width=0.7\columnwidth]{ml/xtend}
		\end{example}
		
	\end{fancycolumns}
\end{frame}


\subsection{Acceleo}
\begin{frame}{\insertsubsection{}}
	\begin{fancycolumns}[animation=none, columns=3, widths={25,45}]
		\begin{example}{Acceleo}
			Code generator compatible with Eclipse Modeling Framework
		\end{example} \pause
		\centering\pic[width=\columnwidth]{ml/acc_input} \pause
		\nextcolumn
		\centering\pic[width=\columnwidth]{ml/acc_template} \pause
		\nextcolumn
		\centering\pic[width=\columnwidth]{ml/acc_output}
	\end{fancycolumns}
\end{frame}

% TODO für die Zukunft
%\subsection{Abstract State Machines}
%\begin{frame}{\insertsubsection{}}
%	\begin{fancycolumns}[animation=none]
%		
%		\nextcolumn
%		
%	\end{fancycolumns}
%\end{frame}

\subsection{Code Completion}
\begin{frame}{\insertsubsection{}}
	\begin{fancycolumns}[animation=none]
		\begin{definition}{Code Completion}
			\begin{itemize}
				\item Suggestion of code lines
				\item Typically during typing at marker characters
				\item Typically as popups
				\item Typically multiple options
				\item Part of almost all modern IDEs
			\end{itemize}
		\end{definition}\pause
		\begin{note}{Goal}
			\begin{itemize}
				\item Improve developer speed
				\item Reduce cognitive load
				\item Prevent errors
			\end{itemize}
		\end{note}\pause
		\nextcolumn
		\begin{note}{Commonly used techniques}
			\begin{itemize}
				\item Static analysis
				\item Adaptive database of source code
				\item User behavior analysis
			\end{itemize}
		\end{note}\pause
		\begin{example}{Code Completion in Visual Studio}
			\centering\pic[width=.9\columnwidth]{ml/code_completion}
		\end{example}
	\end{fancycolumns}
\end{frame}

\subsection{Code Completion with ML}
\begin{frame}{\insertsubsection{}  \mytitlesource{\href{https://dl.acm.org/doi/10.1145/1882362.1882374}{Bruch et al.}}}
		\begin{fancycolumns}[widths={66}]
			\centering\pic[width=\columnwidth]{ml/code_completion_ml}
			\nextcolumn
			\centering\pic[width=\columnwidth]{ml/code_completion_ml_sol}
		\end{fancycolumns}
\end{frame}


%Motivation

%MDSE Model to Code
% MDSE
% Acceleo
% ASM

%Code Completion
% - ML Ansätze
